\chapter{Shortened searches and screens}
\label{ch:abbreviated}

% Author's note (verification and reproduction): session screenlit-q2
% (and screenlit-q2-top); queries q2a-q2e, q2x-q2z; check by meaning
% checks/q2*.json; samples tier2-sample-q2.tsv (seed 20261125) and
% sample60-q2.txt (seed 20261123); see the block at the top of
% search.tex.  Sources: doc/bibliography.bib, keys Ewald2020 Marshall2019
% Hartling2016 Halladay2015 vanEnst2014 Affengruber2020 Pham2016
% Smela2024 Corso2026 Wong2017 Pena2024 Barnish2019 Waffenschmidt2019
% Gartlehner2020 Stoll2019.

A search with a limit per database, a search of fewer databases, a
screen of only part of the records, and a screen by one person instead
of two all save work.
The procedure of \cref{ch:procedure} takes none of the first two
shortcuts and takes the third only with a sample of the rest; for the
fourth it asks for a second screen where the result matters.
This chapter asks what the shortcuts cost when the same reviews are done
both ways: \emph{does an abbreviated search or screen lose few included
studies and leave pooled estimates and conclusions unchanged?}
The claim is tested as a general one; a shortcut that is harmless for
some questions and not for others does not meet it.

\section{Method}
\label{sec:abbreviated-method}

The search is the one described in \cref{ch:lit-search}.
The second question's queries are listed below, and
\cref{tab:abbreviated-matrix} gives their hits per database.
The support queries are q2a to q2d, the counter queries q2x to q2z; q2e
was added after the check by meaning, with the words of methodological
shortcuts, meta-epidemiological comparisons and the value of a second
reviewer.
\begin{description}[style=nextline,font=\normalfont\bfseries]
\item[q2a, support query] {\footnotesize\ttfamily\raggedright (\textquotedbl{}rapid review\textquotedbl{} OR \textquotedbl{}rapid reviews\textquotedbl{} OR abbreviated OR streamlined) AND (\textquotedbl{}systematic review\textquotedbl{} OR \textquotedbl{}systematic reviews\textquotedbl{}) AND (compared OR comparison OR concordance OR agreement)\par}
\item[q2b, support query] {\footnotesize\ttfamily\raggedright (\textquotedbl{}single screening\textquotedbl{} OR \textquotedbl{}single-reviewer screening\textquotedbl{} OR \textquotedbl{}one reviewer\textquotedbl{} OR \textquotedbl{}single reviewer\textquotedbl{}) AND (screening) AND (double OR dual OR \textquotedbl{}two reviewers\textquotedbl{}) AND (\textquotedbl{}systematic review\textquotedbl{})\par}
\item[q2c, support query] {\footnotesize\ttfamily\raggedright (\textquotedbl{}fewer databases\textquotedbl{} OR \textquotedbl{}reduced number of databases\textquotedbl{} OR \textquotedbl{}MEDLINE only\textquotedbl{} OR \textquotedbl{}restricting databases\textquotedbl{} OR \textquotedbl{}limited search\textquotedbl{} OR \textquotedbl{}single database\textquotedbl{}) AND (search) AND (\textquotedbl{}meta-analysis\textquotedbl{} OR \textquotedbl{}systematic review\textquotedbl{}) AND (conclusions OR \textquotedbl{}effect estimates\textquotedbl{} OR \textquotedbl{}pooled estimates\textquotedbl{} OR \textquotedbl{}included studies\textquotedbl{})\par}
\item[q2d, support query] {\footnotesize\ttfamily\raggedright (\textquotedbl{}abbreviated search\textquotedbl{} OR \textquotedbl{}abbreviated searches\textquotedbl{} OR \textquotedbl{}restricted search\textquotedbl{} OR \textquotedbl{}search restrictions\textquotedbl{} OR \textquotedbl{}limited search\textquotedbl{} OR \textquotedbl{}fewer databases\textquotedbl{} OR \textquotedbl{}number of databases\textquotedbl{} OR \textquotedbl{}single database\textquotedbl{}) AND (\textquotedbl{}systematic review\textquotedbl{} OR \textquotedbl{}systematic reviews\textquotedbl{} OR \textquotedbl{}meta-analysis\textquotedbl{} OR \textquotedbl{}meta-analyses\textquotedbl{}) AND (conclusions OR \textquotedbl{}effect estimate\textquotedbl{} OR \textquotedbl{}effect estimates\textquotedbl{} OR \textquotedbl{}pooled estimate\textquotedbl{} OR \textquotedbl{}results changed\textquotedbl{} OR \textquotedbl{}change the results\textquotedbl{})\par}
\item[q2e, support query] {\footnotesize\ttfamily\raggedright (\textquotedbl{}methodological shortcuts\textquotedbl{} OR shortcuts OR \textquotedbl{}rapid review\textquotedbl{} OR \textquotedbl{}rapid reviews\textquotedbl{} OR expedite OR expedited OR \textquotedbl{}second reviewer\textquotedbl{}) AND (\textquotedbl{}systematic review\textquotedbl{} OR \textquotedbl{}systematic reviews\textquotedbl{}) AND (\textquotedbl{}meta-epidemiological\textquotedbl{} OR \textquotedbl{}different results\textquotedbl{} OR \textquotedbl{}case studies\textquotedbl{} OR \textquotedbl{}value of\textquotedbl{} OR \textquotedbl{}same conclusions\textquotedbl{} OR \textquotedbl{}results and conclusions\textquotedbl{})\par}
\item[q2x, counter query] {\footnotesize\ttfamily\raggedright (\textquotedbl{}additional databases\textquotedbl{} OR \textquotedbl{}other databases\textquotedbl{} OR \textquotedbl{}single database\textquotedbl{} OR \textquotedbl{}one database\textquotedbl{}) AND (changed OR altered OR missed OR \textquotedbl{}would have missed\textquotedbl{} OR insufficient) AND (\textquotedbl{}systematic review\textquotedbl{} OR \textquotedbl{}meta-analysis\textquotedbl{})\par}
\item[q2y, counter query] {\footnotesize\ttfamily\raggedright (\textquotedbl{}single screening\textquotedbl{} OR \textquotedbl{}single reviewer\textquotedbl{} OR \textquotedbl{}one reviewer\textquotedbl{}) AND (missed OR missing OR \textquotedbl{}error rate\textquotedbl{} OR \textquotedbl{}false negatives\textquotedbl{}) AND (screening)\par}
\item[q2z, counter query] {\footnotesize\ttfamily\raggedright (\textquotedbl{}rapid review\textquotedbl{} OR \textquotedbl{}rapid reviews\textquotedbl{} OR \textquotedbl{}abbreviated search\textquotedbl{} OR abbreviated OR \textquotedbl{}single screening\textquotedbl{} OR \textquotedbl{}single reviewer\textquotedbl{} OR \textquotedbl{}fewer databases\textquotedbl{}) AND (\textquotedbl{}systematic review\textquotedbl{} OR \textquotedbl{}systematic reviews\textquotedbl{} OR \textquotedbl{}meta-analysis\textquotedbl{}) AND (\textquotedbl{}different conclusions\textquotedbl{} OR \textquotedbl{}changed conclusions\textquotedbl{} OR \textquotedbl{}conclusions changed\textquotedbl{} OR discordant OR \textquotedbl{}missed studies\textquotedbl{} OR \textquotedbl{}missing studies\textquotedbl{} OR \textquotedbl{}would have missed\textquotedbl{})\par}
\end{description}


\begin{table}
  \begin{sidecaption}{%
      Hits per query and database for the second question (shortened
      searches and screens), searched 2026-11-23 with Scholar, as
      \texttt{scholar sessions show} reports them; the databases and
      their syntax as in \cref{tab:stopping-matrix}.
      \enquote{of about}: Semantic Scholar's approximate total;
      \enquote{274 of 275}: the database reported one record more than
      it returned.
    }[tab:abbreviated-matrix]
    \footnotesize
    \begin{tabular}{@{}l*{6}{>{\raggedleft\arraybackslash}p{1.45cm}}@{}}
\toprule
Query & OA & S2 & Scopus & WoS & IEEE & arXiv \\
\midrule
q2a & 907 & 538 of about 541 & 479 & 509 of 510 & 11 & 6 \\
q2b & 332 & 222 of about 226 & 211 & 331 of 332 & 0 & 0 \\
q2c & 322 & 191 of about 192 & 113 & 66 & 1 & 0 \\
q2d & 469 & 336 of about 338 & 274 of 275 & 180 & 3 & 0 \\
q2e & 220 & 321 of about 322 & 165 & 79 & 0 & 1 \\
q2x & 619 & 510 of about 514 & 275 & 343 & 2 & 0 \\
q2y & 197 & 132 & 80 & 73 & 4 & 4 \\
q2z & 42 & 22 & 21 & 16 & 0 & 0 \\
\bottomrule
\end{tabular}

  \end{sidecaption}
\end{table}

Five studies were named before the search (a review of single against
double screening, two reviews of rapid-review methods, and two studies of
what sources beyond PubMed or MEDLINE add); all came back.
The check by meaning returned 36 works no session held; the repair query
returned three of the comparisons cited below and left 33 outside, most
of them general texts on systematic reviews (\cref{sec:search-check}).
Of the 4,123 distinct records, 679 were screened in the first tier, and
38 of them bore on the question (12 supporting the claim, 26 qualifying
or contradicting it).
Of the 3,384 records left unscreened, none among a random sample of 60
bore on the question (95\,\% interval 0 to 6.0\,\%); nor did any of a
second sample of 60 drawn from the records that three or more databases
had returned.
The records judged bearing are listed in \cref{sec:abbreviated-hits}.

\section{Results}
\label{sec:abbreviated-results}

\paragraph{Fewer databases: usually the same answer, not always.}
A meta-epidemiological study redid the searches of 47 Cochrane reviews in
abbreviated forms and found that the abbreviated searches led to
identical effect estimates in 34 to 79\,\% of the meta-analyses and to
opposite effects, or none that could be estimated, in 6 to 13\,\%:
\textcquote[abstract]{Ewald2020}{Abbreviated literature searches often
led to identical or very similar effect estimates as comprehensive
searches with slightly increased confidence intervals. Relevant
deviations may occur.}
In 129 Cochrane reviews from three review groups, restricting the
meta-analyses to the studies two databases held changed little:
\textcquote[abstract]{Hartling2016}{Results of meta-analyses based on the
majority of studies did not differ in most cases. There were very few
cases of changes in statistical significance.}
Studies read only in their abstracts agree: in 50 Cochrane reviews,
\textcquote[abstract]{Halladay2015}{\textins*{m}eta-analysis using only
PubMed-indexed vs.\ all available studies led to a different conclusion
in a single case}; in ten reviews of diagnostic accuracy,
\textcquote[abstract]{vanEnst2014}{\textins*{r}estricting to studies
indexed on MEDLINE did not influence the summary estimates of the
meta-analyses in our sample}; and in three rapid reviews of health
technologies, searching fewer databases would have missed ten references,
yet \textcquote[abstract]{Wong2017}{nine out of the ten references were
found to have no or minimal impact on the overall findings of the
reviews}.
The largest simulation, read only in its abstract, applied rapid-review
shortcuts to 2,512 Cochrane reviews; statistical significance changed in
6.5 to 38.6\,\% of the meta-analyses, with no evidence of bias, and its
authors conclude that \textcquote[abstract]{Marshall2019}{\textins*{s}earching
PubMed only might be considered where a \textasciitilde10\,\% risk of
the primary outcome OR changing by \textgreater20\,\% could be
tolerated}.

\paragraph{Broad questions lose more.}
Abbreviated searches combined with screening by one reviewer gave the
same conclusions as the Cochrane reviews for two oncological topics, but
for a public-health topic the rapid reviews found 38 to 48\,\% of the
included studies, and the authors conclude:
\textcquote[abstract]{Affengruber2020}{For complex public health topics,
sensitivity seems to be insufficient.}
The other comparisons outside focused clinical questions were read only
in their abstracts.
In three reviews in agri-food public health, the shortcuts usually
omitted at least one relevant study, and
\textcquote[abstract]{Pham2016}{\textins*{t}he omission of studies
affected 39 of 143 possible meta-analyses, of which 14 were no longer
possible because of insufficient studies}.
Of three expedited reviews checked against their Cochrane versions, one
missed 3 of 19 studies, and \textcquote[abstract]{Smela2024}{their lack
impacted significantly the final results, or the possibility to run
meta-analyses, in four cases}; a second missed 3 of 27 without changing
the results, and the third matched completely.
Restricting three reviews of low back pain to the three databases the
Cochrane rapid-review guidance names left out studies held by Chinese and
physiotherapy databases: \textcquote[abstract]{Corso2026}{Relative recall
ranged from 65.4\,\% to 100.0\,\%, missing up to 34.6\,\% of relevant
studies.}
In an epidemiological review, a conference abstract reports that
\textcquote[abstract]{Barnish2019}{114 studies were eligible for
inclusion, of which SR identified 101 (89\,\%, due to absence of forward
citation chasing) and RR 64 (56\,\%, McNemar test p\textless0.001)}.
The one comparison found in computing redid a software-engineering
review as a rapid review: its search returned 78\,\% of the papers the
systematic review included, and \textcquote[abstract]{Pena2024}{\textins*{b}oth
secondary studies identified the same number of pieces of evidence (30),
but the pieces of evidence are not the same}.

\section{One screener or two}
\label{sec:abbreviated-single}

A methodological review of single against double screening found that
\textcquote[abstract]{Waffenschmidt2019}{\textins*{t}he median proportion
of missed studies was 5\,\% (range 0 to 58\,\%)}; in 3 of 7 single
screenings whose effect on the findings could be assessed, all by one
less experienced reviewer, the findings of the meta-analyses would have
changed substantially, while the missed studies of the other four had no
or a negligible effect.
In a randomised trial read only in its abstract,
\textcquote[abstract]{Gartlehner2020}{single-reviewer abstract screening
missed 13\,\% of relevant studies (sensitivity: 86.6\,\%; 95\,\%
confidence interval [CI], 80.6\,\%--91.2\,\%)}, against 3\,\% for two
reviewers.
In one large review, also read only in its abstract,
\textcquote[abstract]{Stoll2019}{an additional 6.6\,\% to 9.1\,\% of
eligible studies were identified during the title/abstract stage by using
two reviewers}.
What the rapid-review guidance makes of this, and how it asks for
agreement to be checked, is in \cref{ch:reporting}.

\section{Discussion and limitations}
\label{sec:abbreviated-discussion}

Almost all of this evidence is from health research, and most of it from
reviews of clinical interventions, where the question is focused and a
few databases index the field well.
The comparisons that found larger losses (public health, agri-food,
epidemiology, low back pain, software engineering) were, apart from one,
read only in their abstracts, and so were the largest simulation, the
two studies of sources beyond PubMed and MEDLINE, the trial and the
single review on second reviewers: their full texts were behind bot
checks, refused by the publisher, or not tried.
No source tests a limit on the number of results per database as such,
which is the shortcut that \cref{sec:measured-limit} measured; the
nearest is the study of a ranked web source in \cref{ch:stopping}.

\section{Conclusion}
\label{sec:abbreviated-conclusion}

Not in general: for focused clinical questions, searching fewer
databases usually left pooled estimates identical or close, but changed
them in a minority of meta-analyses; for broad questions and outside
clinical medicine the shortcuts lost a large share of the included
studies; and one screener missed more relevant studies than two.
For the procedure this means searching every live database without a
limit, saying what was not screened, and, where a result matters,
screening a second time.

\section{Future work}
\label{sec:abbreviated-future}

This search left eleven of the fifteen sources read only in their
abstracts; their full texts would give the conditions of each loss.
The field lacks a comparison of the full and the shortened search in
computing and education beyond the one software-engineering case, and a
test of a limit per database on what a review concludes; if such a test
found that a limit rarely changes conclusions in these fields, the
procedure could spend less on searching, and if not, the case for no
limit would rest on more than one review.

\section{The records judged bearing}
\label{sec:abbreviated-hits}

% Author's note: table by generate_audit_table on a copy of screenlit-q2,
% bearing rows only, cited = cited in doc/bibliography.bib.
% --- scholar LaTeX fragment --------------------------------------------
% Generated by: scholar sessions export -f table
% Session: screenlit-q2
% \input this file from the body of a document whose
% preamble loads:
%   \usepackage{longtable}
%   \usepackage{booktabs}
%   \usepackage{array}
% ----------------------------------------------------------------------

\begingroup\footnotesize
\begin{longtable}{@{}%
  >{\raggedright\arraybackslash}p{0.44\textwidth}c%
  >{\raggedright\arraybackslash}p{0.13\textwidth}%
  >{\raggedright\arraybackslash}p{0.26\textwidth}@{}}
\caption{Records bearing on the claim, Question 2, abbreviated searches and screens (16 cited, 7 supporting and 17 qualifying or contradicting it, out of 3,922 unique hits; the 82 adjacent hits, the 670 false hits and 3,130 unclassified are not listed). The last column gives the reason for each inclusion; for rows decided by an agent the agent's confidence is given. A row is cited when the text cites the work, whoever decided on its records. Databases searched: OpenAlex (OA; not open access), IEEE Xplore (IEEE), Web of Science (WoS), Scopus, Semantic Scholar (S2), arXiv. Searched 2026-11-23.}\label{tab:screenlit-q2}\\
\toprule
Title & Year & Database & Reason \\
\midrule
\endfirsthead
\multicolumn{4}{@{}l}{\emph{\tablename~\thetable{} (cont.)}}\\
\toprule
Title & Year & Database & Reason \\
\midrule
\endhead
\midrule \multicolumn{4}{r@{}}{\emph{continued on next page}}\\
\endfoot
\bottomrule
\endlastfoot
Abbreviated and comprehensive literature searches led to identical or very similar effect estimates: a meta-epidemiological study & 2020 & WoS, Scopus, OA, S2 & \emph{cited}: supports the claim (0.80) \\
An evaluation of DistillerSR's machine learning-based prioritization tool for title/abstract screening -- impact on reviewer-relevant outcomes & 2020 & OA, WoS, Scopus, S2 & \emph{cited}, pending, unclassified \\
Assessing the coverage of PubMed, Embase, OpenAlex, and Semantic Scholar for automated single-database searches in living guideline evidence surveillance: a case study of the international polycystic ovary syndrome guidelines 2023 & 2025 & Scopus & \emph{cited}, pending, unclassified \\
Assessing the validity of a rapid review against a systematic literature review. A comparison of systematic literature reviews done by Cochrane with rapid reviews and the impact on meta-analyses results & 2024 & OA, WoS, Scopus, S2 & \emph{cited}: qualifies or contradicts the claim (0.70) \\
Combining abbreviated literature searches with single-reviewer screening: three case studies of rapid reviews & 2020 & OA, WoS, Scopus, S2 & \emph{cited}: qualifies or contradicts the claim (0.75) \\
Comparing the Efficacy of Rapid Review With a Systematic Review in the Software Engineering Field & 2024 & OA, WoS, Scopus, S2 & \emph{cited}: qualifies or contradicts the claim (0.70) \\
Does a rapid review version of a large epidemiological systematic review fail to identify many eligible studies, and what implications does this have for the results of the review? & 2019 & OA, S2 & \emph{cited}: qualifies or contradicts the claim (0.60) \\
Evaluating the proportion of retrieved citations and relevant studies when database selection is restricted to MEDLINE, CENTRAL, and Embase: part 4 of a meta-research study of rapid review methodological trade-offs & 2026 & OA, WoS, Scopus, S2 & \emph{cited}: qualifies or contradicts the claim (0.70) \\
Implications of applying methodological shortcuts to expedite systematic reviews: three case studies using systematic reviews from agrifood public health & 2016 & OA, S2, WoS, Scopus & \emph{cited}: qualifies or contradicts the claim (0.75) \\
Meta-epidemiologic analysis indicates that MEDLINE searches are sufficient for diagnostic test accuracy systematic reviews & 2014 & OA, WoS, Scopus, S2 & \emph{cited}: supports the claim (0.70) \\
Rapid reviews may produce different results to systematic reviews: a meta-epidemiological study & 2018 & OA, S2, WoS, Scopus & \emph{cited}: qualifies or contradicts the claim (0.80) \\
Single-reviewer abstract screening missed 13 percent of relevant studies: a crowd-based, randomized controlled trial & 2020 & OA, WoS, Scopus, S2 & \emph{cited}: qualifies or contradicts the claim (0.75) \\
The contribution of databases to the results of systematic reviews: a cross-sectional study & 2016 & OA, WoS, Scopus, S2 & \emph{cited}: supports the claim (0.70) \\
The Impact Of Searching Fewer Databases In Health Technology Assessment Rapid Reviews & 2017 & OA, S2 & \emph{cited}: supports the claim (0.65) \\
The value of a second reviewer for study selection in systematic reviews & 2019 & OA, S2, WoS, Scopus & \emph{cited}: qualifies or contradicts the claim (0.70) \\
Using data sources beyond PubMed has a modest impact on the results of systematic reviews of therapeutic interventions & 2015 & Scopus & \emph{cited}: supports the claim (0.75) \\
Empirical Testing of Alternative Search Methods to Retrieve Utility Values for Health Economic Modelling & 2024 & OA, WoS, Scopus, S2 & supports the claim (0.50) \\
Excluding non-English publications from evidence-syntheses did not change conclusions: a meta-epidemiological study & 2019 & OA, S2, WoS, Scopus & supports the claim (0.60) \\
Increasing the efficiency of study selection for systematic reviews using prioritization tools and a single-screening approach & 2023 & OA, Scopus, WoS, S2 & supports the claim (0.55) \\
Limited search strategies were effective in finding relevant nonrandomized studies & 2006 & Scopus, S2, OA, WoS & supports the claim (0.50) \\
Limits in the search date for rapid reviews of diagnostic test accuracy studies & 2022 & OA, WoS, Scopus, S2 & supports the claim (0.60) \\
Simulated changes in systematic review screening sensitivity and specificity associated with varying co-screening proportions & 2026 & OA, Scopus, S2 & supports the claim (0.55) \\
Successful incorporation of single reviewer assessments during systematic review screening: development and validation of sensitivity and work-saved of an algorithm that considers exclusion criteria and count & 2021 & OA, WoS, Scopus, S2 & supports the claim (0.55) \\
A retrospective comparison of systematic reviews with same-topic rapid reviews & 2017 & OA, WoS, Scopus, S2 & qualifies or contradicts the claim (0.60) \\
Characteristics and recovery methods of studies falsely excluded during literature screening---a systematic review & 2022 & OA, WoS, Scopus, S2 & qualifies or contradicts the claim (0.45) \\
Cochrane systematic reviews in acupuncture: Methodological diversity in database searching & 2005 & WoS, OA, S2, Scopus & qualifies or contradicts the claim (0.45) \\
Comparison of a full systematic review versus rapid review approaches to assess a newborn screening test for tyrosinemia type 1 & 2017 & OA, WoS, Scopus, S2 & qualifies or contradicts the claim (0.60) \\
Database selection in systematic reviews: an insight through clinical neurology & 2017 & S2, OA, WoS, Scopus & qualifies or contradicts the claim (0.50) \\
Error rates of human reviewers during abstract screening in systematic reviews & 2020 & OA, WoS, Scopus, S2 & qualifies or contradicts the claim (0.55) \\
Evaluating the accuracy and impact of methodological trade-offs used in rapid reviews on results and certainty of the evidence: meta-research studies in systematic reviews of chronic low back pain & 2026 & OA & qualifies or contradicts the claim (0.65) \\
Evaluating the accuracy of full text and serial screening by one vs two independent reviewers: part 2 of a meta-research study of rapid review methodological trade-offs & 2026 & OA, WoS, Scopus, S2 & qualifies or contradicts the claim (0.60) \\
Evaluating the accuracy of screening titles and abstracts for progressing to full-text screening by one vs two independent reviewers: part 1 of a meta-research study of rapid review methodological trade-offs & 2026 & OA, WoS, Scopus, S2 & qualifies or contradicts the claim (0.65) \\
Evaluating the impact of methodological trade-offs on meta-analyses and certainty of the evidence according to GRADE: part 5 of a meta-research study of rapid reviews & 2026 & OA, WoS, Scopus, S2 & qualifies or contradicts the claim (0.60) \\
Examining screening test accuracy in Tyrosinemia type 1: A comparison of rapid and systematic review Methods & 2016 & OA, S2 & qualifies or contradicts the claim (0.50) \\
Identification of randomized controlled trials in systematic reviews: accuracy and reliability of screening records & 2002 & OA, WoS, Scopus, S2 & qualifies or contradicts the claim (0.60) \\
Lessons for search strategies from a systematic review, in The Cochrane Library, of nutritional supplementation trials in patients after hip fracture & 2001 & WoS, OA, Scopus, S2 & qualifies or contradicts the claim (0.55) \\
Literature search in systematic reviews: How much is good enough? & 2023 & OA, WoS, Scopus & qualifies or contradicts the claim (0.60) \\
Literature searching for social science systematic reviews: consideration of a range of search techniques & 2010 & OA, WoS, Scopus, S2 & qualifies or contradicts the claim (0.50) \\
Network meta-analyses could be improved by searching more sources and by involving a librarian & 2014 & OA, S2, WoS, Scopus & qualifies or contradicts the claim (0.45) \\
Sources of evidence for systematic reviews of interventions in diabetes & 2005 & WoS, OA, Scopus & qualifies or contradicts the claim (0.55) \\
\end{longtable}
\endgroup

